\begin{table}[htbp]\centering
\begin{threeparttable}
\caption{Assignment of the sequences of withheld cultures and species, with the external databases intact.}
\label{tab:benchmark}
\footnotesize\setlength{\tabcolsep}{4pt}
\begin{tabular}{lrrrr}
\toprule
LSU naming & Known kept & Exact species & False known & False known\\
 & (of 755) & & species withheld & single-culture\\
 & & & (of 803) & species (of 48)\\
\midrule
Nearest LSU reference & 748 & 603 & 215 (26.8\%) & 3\\
LSU species clade (placement) & 751 & 701 & 71 (8.8\%) & 0\\
Placement + ITS distance rule & 745 & 697 & 18 (2.2\%) & 0\\
\bottomrule
\end{tabular}
\begin{tablenotes}\footnotesize
\item \emph{Known kept}: main-type sequences of species represented by another remaining culture that were
assigned to a known species or species group that includes their own species (148 rounds, one culture withheld in
each). \emph{False known}: sequences assigned to a known species although their species has no remaining culture,
in 50 rounds with one species withheld in each and, in the rounds with one culture withheld, among the main-type
sequences of species with a single culture. Averaged over species and over cultures, respectively, the final
false-known rate with a species withheld is 1.4\% and 1.9\%.
\end{tablenotes}
\end{threeparttable}
\end{table}
